\documentclass{article}
\usepackage{PRIMEarxiv} % Core package for the PrimeAI Template
\usepackage{amsmath, amssymb, amsthm, bm, dcolumn} % Add amsthm here for the proof environment
\usepackage[numbers,sort&compress]{natbib} % Natbib for citations
\usepackage{graphicx} % For high-quality images
\usepackage{multicol} % Multi-column figures/tables
\usepackage{hyperref} % Hyperlinks
\usepackage{listings} % Code listings
\usepackage{authblk} % For structured affiliations
% Define theorem style (optional)
\theoremstyle{plain}
\newtheorem{theorem}{Theorem}
\renewcommand{\qedsymbol}{}

% Custom commands for primes and qubit encoding
\newcommand{\primeQubit}[1]{|p_{#1}\rangle}
\newcommand{\primeGate}[1]{U_{p_{#1}}}

\usepackage{wrapfig}
\usepackage[pscoord]{eso-pic}
\usepackage[fulladjust]{marginnote}
\reversemarginpar

% Typesetting improvements without footnote patching
\usepackage[protrusion=true, expansion=true, tracking=false]{microtype}
\microtypecontext{spacing=nonfrench}

% Line numbers
\usepackage[right]{lineno}

% Text layout - adjust as needed
\raggedright
\setlength{\parindent}{0.5cm}
\textwidth 5.25in 
\textheight 8.75in

% Set double spacing
\usepackage{setspace} 
\doublespacing

% Adjust width for specific content
\usepackage{changepage}

% Adjust caption style
\usepackage[aboveskip=1pt,labelfont=bf,labelsep=period,singlelinecheck=off]{caption}

% Remove brackets from references
\makeatletter
\renewcommand{\@biblabel}[1]{\quad#1.}
\makeatother

% Header, footer, and page numbers
\usepackage{lastpage,fancyhdr}
\pagestyle{fancy}
\fancyhf{}
\fancyhead[L]{Preprint - PrimeAI Enhanced Template}
\fancyfoot[C]{\scriptsize Multiplicity Theory © 2024 Ryan Van Gelder - Citizen Gardens \\ Licensed Under MIT and CC BY-NC-SA 4.0.}
\fancyfoot[R]{Page \thepage\ of \pageref{LastPage}}
\renewcommand{\footrule}{\hrule height 2pt \vspace{2mm}}

\begin{document}
\title{Quantum Musicality: Applying Multiplicity Theory to Time-Evolving Musical Compositions}

\author{Ryan Van Gelder}
\affil{Citizen Gardens - The Foundation of Multiplicity}
\date{\today}
\maketitle

\begin{abstract}
Quantum Musicality introduces a novel approach to composing time-evolving musical pieces by integrating principles from Multiplicity Theory. By leveraging harmonic resonance, prime-based encoding, tensor networks, and quantum coherence, this framework facilitates dynamic, adaptable, and multidimensional musical creations. Applications span generative systems, quantum-inspired instruments, and adaptive compositions, providing groundbreaking opportunities for interdisciplinary exploration.
\end{abstract}

\section{Introduction}
Multiplicity Theory provides a robust framework for understanding interconnected systems through harmonic interactions, recursive feedback, and modular encoding. In this paper, we extend its principles to music, conceptualizing quantum musicality as a paradigm where compositions evolve dynamically, reflecting quantum principles like superposition, entanglement, and coherence.

\section{Foundational Principles of Quantum Musicality}

\subsection{Multiplicity and Harmonic Resonance}
The dynamics of harmonic resonance can be modeled using the Enhanced Dynamic Multiplicity Equation \cite{galioto2024dynamic}:
\begin{equation}
\frac{\partial \rho_k}{\partial t} = \alpha_k(t) \rho_k + \beta_k(t) I_k + \gamma_k(t) \sum_j T_{kj} \rho_j + \lambda(t) \big(\Omega_B(\rho) + \Omega_{FS}(\rho)\big),
\end{equation}
where $\rho_k$ represents the state of a musical component, and $\alpha_k$, $\beta_k$, $\gamma_k$, and $\lambda(t)$ are time-dependent parameters modulating interactions among tones.

\subsection{Prime-Based Encoding}
Musical tones, intervals, and rhythms can be encoded using prime numbers:
\begin{equation}
\phi(v_i) = p_k, \quad \phi(e_i) = \prod_{j \in e_i} p_j,
\end{equation}
where $v_i$ represents a tone, $e_i$ an interval, and $p_k$ a prime assigned to each entity. This encoding enables hierarchical and modular composition.

\subsection{Tensor Networks and Polyphonic Structures}
Tensor networks are used to represent polyphonic and harmonic dependencies:\
\begin{equation}
T(t) = \sum_{i,j,k} T_{ijk} \otimes \phi(p_{ijk}),
\end{equation}
where $T_{ijk}$ encodes the interaction of multiple musical elements, and $\phi(p_{ijk})$ represents their prime-based mappings.

\subsection{Quantum Coherence in Music}
Quantum coherence can model overlapping musical states:\
\begin{equation}
|\psi(t)\rangle = \sum_{i=1}^N \alpha_i(t) |\phi_i\rangle,
\end{equation}
where $|\phi_i\rangle$ represents a musical state and $\alpha_i(t)$ its time-evolving amplitude.

\section{Applications and Framework}

\subsection{Interactive Musical Systems}
Leveraging **recursive feedback** mechanisms, interactive systems can generate real-time adaptive compositions:\
\begin{equation}
M(t+1) = f\big(M(t), R(t)\big),
\end{equation}
where $M(t)$ is the current state of the system, and $R(t)$ provides recursive feedback from live input.

\subsection{Quantum-Inspired Instruments}
Quantum principles inform instrument design, enabling musicians to manipulate quantum states. For example, the evolution of sound can be described using a stochastic term:
\begin{equation}
P(S, t) = \prod_{i=1}^N \big(p(x_i, t) + \epsilon_i(t)\big),
\end{equation}
where $\epsilon_i(t)$ represents noise or fluctuation in the generated tone.

\subsection{Compositional Algorithms}
Generative algorithms can utilize the Dynamic Multiplicity Equation:\
\begin{equation}
\frac{\partial \rho_k}{\partial t} = \sum_{j} \gamma_{kj} \rho_j \cos(\theta_k(t) - \theta_j(t)),
\end{equation}
where $\theta_k(t)$ represents the phase evolution of $\rho_k$.

\section{Future Directions}

\subsection{Simulation Tools}
To simulate harmonic interactions within tensor networks, consider representing the system's state as a high-dimensional tensor:
\[
\mathcal{T}(t) = \sum_{i,j,k} T_{ijk}(t) \, \phi(p_{ijk}),
\]
where $T_{ijk}(t)$ represents the interaction coefficients evolving over time, and $\phi(p_{ijk})$ encodes the prime-based structure of the musical elements. The simulation can be iteratively updated using:
\[
\mathcal{T}(t+1) = f\big(\mathcal{T}(t), R(t)\big),
\]
where $R(t)$ is a feedback term modeling environmental or user interactions.

\subsection{Encoding Frameworks}
The prime-based encoding framework can be enhanced to include dynamic phase shifts for polyphonic integration:
\[
\phi_k = e^{2\pi i n_k / p_k} \cdot e^{i \theta_k(t)},
\]
where $p_k$ is a prime assigned to a musical component, $n_k$ is the component's discrete state, and $\theta_k(t)$ represents the phase shift modulated over time. The encoded polyphonic structure can be expressed as:
\[
\Phi(t) = \prod_{k=1}^N \phi_k(t).
\]

\subsection{Hybrid Computing}
Hybrid quantum-classical systems can leverage the strengths of both paradigms to manage large-scale simulations. The state of the hybrid system can be represented as:
\[
| \Psi(t) \rangle = \sum_{k=1}^N \alpha_k(t) | q_k \rangle \otimes \mathbf{c}_k,
\]
where $| q_k \rangle$ represents the quantum states, $\mathbf{c}_k$ are the classical computational states, and $\alpha_k(t)$ are time-evolving coefficients. The evolution can be governed by:
\[
\frac{d | \Psi(t) \rangle}{dt} = \mathcal{H}(t) | \Psi(t) \rangle + \mathcal{C}(t),
\]
where $\mathcal{H}(t)$ is the quantum Hamiltonian and $\mathcal{C}(t)$ incorporates classical interactions.

\subsection{Generative AI Models}
AI-driven generative models can incorporate multiplicity principles through recursive feedback and hierarchical representations. Let the generative function $G$ for a composition be defined as:
\[
G(\mathbf{z}, t) = \sum_{k=1}^N \phi_k(t) \cdot h_k(\mathbf{z}),
\]
where $\mathbf{z}$ is a latent vector, $h_k$ represents harmonic components influenced by multiplicity principles, and $\phi_k(t)$ are the prime-based encodings. Recursive adaptation can be implemented via:
\[
\mathbf{z}(t+1) = f(\mathbf{z}(t), \nabla G(\mathbf{z}(t), t)),
\]
where $f$ adjusts the latent variables based on the gradient of the generative function.


\section{Conclusion}
Quantum Musicality represents a bold step forward in integrating Multiplicity Theory into the realm of music. By combining harmonic principles, quantum coherence, and advanced encoding, this paradigm opens the door to revolutionary applications in adaptive music systems, quantum instruments, and beyond.


\nocite{*}
\bibliographystyle{plain}
\bibliography{references}

\end{document}